\section{The Bulk as a Data Structure}

Earlier sections built the mesh and ran physics on it. This one reads the same mesh as a machine for holding and finding data. The claim is sharp. The hyperbolic bulk is not a place where you might bolt a database. It already is one. A real storage engine sorts itself into the crystal region by region, and each region holds the part of the engine its own geometry is built for.

The mesh is the whole $\{3,4,3,4\}$ hyperbolic honeycomb. A dock is one cell of that honeycomb, a here-now. A site is one of the $24$ directions out of a dock, and the wake is the growing edge where new docks are born. These four atoms, mesh, dock, site, and wake, are all we need. We never add a stored pointer, a routing table, or an index on the side. The geometry supplies them for free.

The throughline is one sentence. The tree-shaped structures, which are most of a database, live in the bulk with no stored pointers and automatic balance. The array-shaped structures live on the flat cusp. The hash-shaped structures live on the exponential boundary. The append log runs along the wake. Every result below was established on the actual generated mesh, each with a flat control, across twenty-five experiments.

\begin{result}[The bulk is already a database engine]
A hyperbolic crystal is not a place to attach a database to. It is structurally already a storage and computing medium, each region holding the part of the engine its geometry is built for. The parts it is bad at, dense arrays and contiguous range scans, are exactly the parts that belong on the flat cusp, not in the bulk.
\end{result}

\subsection{The bulk as a computer, read as memory}

A companion section showed that the knit, the collide-then-stream law of the mesh, can compute anything. The knit is universal. This section answers a different question. It shows where the data goes. The same crystal that runs the universal flow is, read as memory, a balanced index with logarithmic depth.

Two facts do all the work, and everything else is a consequence of them.

\begin{table}[ht]
\centering
\begin{tabular}{@{}ll@{}}
\toprule
fact & consequence \\
\midrule
logarithmic depth & a ball of radius $R$ holds the whole structure, every descent is $O(\log N)$ \\
exponential capacity & a ball of radius $R$ holds about $e^{(n-1)R}$ docks \\
\bottomrule
\end{tabular}
\caption{The two facts behind the bulk as memory. Here $n$ is the bulk dimension and $N$ the number of docks.}
\label{tab:two-facts}
\end{table}

Capacity is exponential in radius. So a unique address has length only logarithmic in the dock count. So every tree descent, every distributed-hash route, and every Merkle proof finishes in $O(\log N)$ steps. The negative curvature of the bulk is the reason. It makes greedy routing and tree embedding succeed where flat space fails.

\begin{principle}[Exponential capacity forces logarithmic everything]
Because the number of docks within radius $R$ grows like $e^{(n-1)R}$, any one dock is named by an address of length $O(R) = O(\log N)$. Every operation that walks from a dock to the center, or between two docks through the center, therefore costs $O(\log N)$.
\end{principle}

\subsection{The six primitives that make the bulk a data structure}

Six geometric facts are the levers. Each is a property the bulk has on its own, and each turns into a data-structure capability with nothing added.

\begin{table}[ht]
\centering
\begin{tabular}{@{}ll@{}}
\toprule
primitive & the fact \\
\midrule
exponential capacity & a ball of radius $R$ holds about $e^{(n-1)R}$ docks \\
logarithmic addressing & a unique canonical address of length $O(\log N)$ per dock \\
implicit neighbours & the dock is its integer $D_4$ coordinate, neighbours are arithmetic \\
tree-likeness & greedy routing succeeds, trees embed at distortion near $1$ \\
the radial scale axis & the Busemann coordinate is a multiresolution pyramid \\
the bulk shortcut & the diameter is logarithmic \\
\bottomrule
\end{tabular}
\caption{The six primitives. The radial axis is the scale. The diameter is logarithmic.}
\label{tab:six-primitives}
\end{table}

The third primitive deserves a word, because it is the one that erases the cost of adjacency. Each dock is its own integer $D_4$ coordinate. Its $24$ neighbours are obtained by arithmetic on that coordinate, one per site. There is no adjacency list stored anywhere. The dock knows who its neighbours are because it knows where it is. This is why the structures below carry zero stored child pointers. The pointer is the geometry.

\subsection{Greedy routing in the bulk}

The headline navigation result is that in the bulk you can route with no routing table at all.

Each dock knows only two things, its own coordinate and the coordinates of its $24$ neighbours. To send a message to a target, a dock forwards it to whichever neighbour is closest to the target in the hyperbolic metric, then that neighbour does the same. Repeat until you arrive. No dock holds a map of the mesh.

\begin{table}[ht]
\centering
\begin{tabular}{@{}ll@{}}
\toprule
result & measurement \\
\midrule
greedy delivery on the geometric embedding & $100$ percent \\
greedy delivery scrambled & $3$ percent \\
tree embedding distortion in the disk & $1.3$ \\
tree embedding distortion in the plane & unbounded \\
\bottomrule
\end{tabular}
\caption{Greedy routing and tree embedding, measured on the generated mesh against a flat control.}
\label{tab:greedy}
\end{table}

\begin{result}[Greedy routing always arrives]
Greedy routing succeeds because negative curvature makes the straight path through the bulk the short path. In flat space the same greedy walk gets stuck in local minima. The bulk has no long thin paths, so there are no traps.
\end{result}

This is the geometric reason a tree built on the bulk cannot degenerate into a list. A list is the worst case of a tree, a long thin path. The bulk has none. The scrambled control, where the same docks are connected without respecting the metric, falls from full delivery to near zero, which shows the success is geometric and not an artifact of the connectivity.

\subsection{The bulk holds the tree-shaped structures with no stored pointers}

This is the central result. The chamber tree of the tessellation is the database tree. Its children are physical neighbour docks. It cannot degenerate, because the geometry has no long thin paths to degenerate along.

\begin{table}[ht]
\centering
\begin{tabular}{@{}ll@{}}
\toprule
structure & result \\
\midrule
B-tree, the index & a point query is a $3$-step physical descent for $4000$ docks, zero stored child pointers \\
trie, radix & every address is its parent address plus one digit, prefixes are physical \\
heap & the radial depth is a heap order, peek-min is the root in $O(1)$ \\
Merkle DAG & an inclusion proof is a logarithmic-length path to the root \\
union-find & find depth is logarithmic from the short bulk paths \\
distributed hash table & a lookup routes up to the common prefix and down in $O(\log N)$ hops, $O(1)$ state \\
\bottomrule
\end{tabular}
\caption{Tree-shaped structures in the bulk, each instantiated with no stored child pointers.}
\label{tab:tree-structures}
\end{table}

\begin{theorem}[The B-tree is the chamber tree]
A B-tree of order $24$ instantiates on $\{3,4,3,4\}$ with no stored child pointers and answers a point query in a logarithmic-depth physical descent. The most-used database structure is the cheapest thing the bulk can hold.
\end{theorem}

The B-tree order is exactly the dock coordination, $24$, the same $24$ sites out of every dock. The fan-out of the index is the fan-out of space. A point query over four thousand docks resolves in a three-step physical walk from one dock to a neighbour to a neighbour, each step chosen by arithmetic on coordinates. No node stores where its children are. The children are simply the adjacent docks.

\subsection{The boundary holds the hash-shaped structures}

The exponential boundary of the bulk is a huge, cheap address space. There are far more docks near the boundary than anywhere inside, so the boundary is where structures that want a vast spread of distinct addresses belong.

\begin{table}[ht]
\centering
\begin{tabular}{@{}ll@{}}
\toprule
structure & result \\
\midrule
hash table & keys hash to exact integer addresses, $O(1)$ probes at half load \\
Bloom filter & the false-positive rate falls from $1.0$ to $0.04$ as the boundary grows \\
inverted index & terms hash to boundary docks, retrieval is output-sensitive \\
\bottomrule
\end{tabular}
\caption{Hash-shaped structures on the exponential boundary.}
\label{tab:hash-structures}
\end{table}

A hash table here is not an approximation of a hash table. The integer $D_4$ coordinate of a boundary dock is the exact address a key hashes to. The probe is reading that dock. At half load the probe count is constant, the standard hash-table guarantee, delivered by the geometry of the boundary rather than by an array allocated in advance.

\subsection{The radial axis holds the scale-and-level structures}

The Busemann coordinate measures how deep a dock sits along the radial axis, from the center outward. It is a clean scale axis. Structures whose job is to organize data by level live along it.

\begin{table}[ht]
\centering
\begin{tabular}{@{}ll@{}}
\toprule
structure & result \\
\midrule
LSM-tree & the radial shells are geometric sorted levels, logarithmic level count \\
multiresolution mipmap & the Busemann levels form a pyramid, fine levels hold geometrically more docks \\
R-tree & nested horoballs are a bounding-volume hierarchy, logarithmic levels \\
skip list & the diameter is logarithmic, far below the flat square-root radius \\
\bottomrule
\end{tabular}
\caption{Scale-and-level structures along the radial axis.}
\label{tab:radial-structures}
\end{table}

Each radial shell holds geometrically more docks than the one inside it, by the exponential capacity. So the shells are already sorted into levels of geometrically growing size, which is exactly the shape of a log-structured merge tree and of an image mipmap. The number of levels is logarithmic in the total dock count, because the capacity is exponential in the radius.

\subsection{Capacity, embedding, traversal, and order}

A further family of standard structures lands directly, covering associative recall, embedding fidelity, traversal, and total order.

\begin{table}[ht]
\centering
\begin{tabular}{@{}ll@{}}
\toprule
structure & result \\
\midrule
associative memory & binding capacity scales with dock count, $0.95$ versus $0.28$ recall \\
tree embedding & distortion $1.3$ in the disk, unbounded in the plane \\
greedy routing & $100$ percent delivery on the geometric embedding, $3$ percent scrambled \\
sorting & the canonical address is a total order, sorting is reading in address order \\
BFS, DFS & the BFS frontier is the wake shell, traversal needs no queue \\
list, stack, queue & a list is a dock path, a stack is a radial ray, $O(1)$ per step \\
\bottomrule
\end{tabular}
\caption{Capacity, embedding, traversal, and order on the bulk.}
\label{tab:capacity-structures}
\end{table}

Two rows are worth slowing down on. Sorting needs no comparison sweep, because the canonical address is already a total order. To sort is to read the docks in address order. And breadth-first search needs no queue, because the frontier of the search is the wake shell itself, the growing edge of the mesh. The traversal data structure is the geometry of the frontier. The associative-memory row, where recall holds at $0.95$ against a flat-control $0.28$, is the seed of the associative-memory engine taken up in a later section, where the same exponential capacity is measured in depth.

\subsection{The honest negatives, what the bulk is bad at}

The methodology requires reporting the costs, not only the wins. The same exponential growth that gives capacity also costs three things, and they are real.

\begin{table}[ht]
\centering
\begin{tabular}{@{}ll@{}}
\toprule
caveat & result \\
\midrule
the interior is nearly empty & the boundary holds the dominant fraction of docks \\
range scans are exponential & a radius-$\rho$ range visits exponentially many docks \\
the bulk is not an array & the interior is too sparse, a dense array belongs on the flat cusp \\
\bottomrule
\end{tabular}
\caption{The honest negatives. What the bulk does badly.}
\label{tab:negatives}
\end{table}

The bulk is excellent for hierarchies, routing, fault-tolerant storage, sketches, and multiresolution. It is poor for dense uniform arrays and contiguous range scans. A range scan asks for every dock within some radius, and the bulk answers with exponentially many, because that is how many there are. A dense array wants uniform spacing, and the bulk interior is too sparse for it. These are not flaws to patch. They are the geometry telling you the dense-array part of the engine lives elsewhere, on the flat cusp.

\subsection{The dimension scaling}

The same structures sharpen or soften with the bulk dimension, by one rule. The capacity exponent equals the dimension minus one.

\begin{table}[ht]
\centering
\begin{tabular}{@{}llll@{}}
\toprule
property & 2D ($\{7,3\}$, $\{5,4\}$) & 3D ($\{5,3,4\}$) & 4D ($\{3,4,3,4\}$) \\
\midrule
capacity in radius $R$ & $e^{R}$ & $e^{2R}$ & $e^{3R}$ \\
flat data it stores & sequences & images, grids & volumes \\
fan-out & $3$ to $7$ & $12$ & $24$ \\
native home for & embeddings, routing, the cleanest tries & image mipmaps, the physics-co-located store & wide database indexes, exact-address hash tables \\
\bottomrule
\end{tabular}
\caption{Dimension scaling. The capacity exponent is the dimension minus one.}
\label{tab:dimension-scaling}
\end{table}

The four-dimensional substrate of this paper, $\{3,4,3,4\}$, sits at the high end. Its fan-out of $24$ and its $e^{3R}$ capacity make it the native home for wide database indexes and exact-address hash tables. A later section extends this table across all forty-two buildable tessellations from two to five dimensions, but the rule never changes. The exponent is the dimension minus one.

\subsection{The bottom line}

Take the dictionary that maps one world onto the other. A dock is a slot. A site, the directed edge out of a dock, is a pointer. A walk along sites is an operation. An address is a location. Under that dictionary a real storage engine sorts itself into the crystal.

\begin{itemize}
\item the B-tree index is the bulk chamber tree, balanced by geometry with zero stored pointers
\item the hash indexes are on the boundary
\item the LSM levels are the radial shells
\item the column data is on the cusp
\item the write-ahead log is the wake
\end{itemize}

\begin{result}[The structural identity]
A hyperbolic crystal is structurally already a database engine, each region holding the part of the engine its geometry is built for. The parts it is bad at, dense arrays and range scans, are exactly the parts that belong on the flat cusp, not in the bulk.
\end{result}

All of this rests on twenty-five experiments, every one passing with a flat control, covering the data structures in the bulk and the standard structures on tessellations. The associative-memory engine, the universal-computation section, and the survey of data structures across all buildable tessellations carry the story further. The same boundary read as a screen is the subject of holography. The mesh stores, the knit computes, and the wake logs, all from the same eight atoms and nothing added.
